\documentclass{beamer}

% Theme and color
\usetheme{Madrid}
\usecolortheme{default}

\definecolor{mnitblue}{RGB}{54,51,180}

\setbeamercolor{structure}{fg=mnitblue}
\setbeamercolor{frametitle}{bg=mnitblue, fg=white}
\setbeamercolor{title}{bg=mnitblue, fg=white}

% Presentation details
\title{Introduction to Beamer}
\author{Siddarth Kumar}
\institute{MNIT Jaipur}
\date{\today}

\begin{document}

\begin{frame}
    \titlepage
\end{frame}

\begin{frame}{Introduction}
    \textbf{What is Beamer?}

    \begin{itemize}
        \item Beamer is a LaTeX document class used to create presentations.
        \item It allows us to create structured slides.
        \item It supports equations, tables, figures and more.
    \end{itemize}
\end{frame}
\begin{frame}{Document and Page Styles}

    \textbf{Document Class}

    \begin{itemize}
        \item Defines the overall structure of a document.
        \item Examples:
        \begin{itemize}
            \item article
            \item report
            \item book
        \end{itemize}
    \end{itemize}

    \vspace{0.3cm}

    \textbf{Page Style}

    \begin{itemize}
        \item Controls the appearance of headers and footers.
        \item Examples: plain, empty, headings and fancy.
    \end{itemize}

\end{frame}


\begin{frame}{Mathematical Equations}

    Beamer also supports mathematical expressions.

    \[
        E = mc^2
    \]

    \[
        \sum_{i=1}^{n} i = \frac{n(n+1)}{2}
    \]

    \textbf{Advantages}
    \begin{itemize}
        \item Mathematical expressions are easy to format.
        \item Equations can be displayed clearly on slides.
    \end{itemize}

\end{frame}
\begin{frame}{Conclusion}

    \begin{center}
        \Huge \textbf{Thank You!}

        \vspace{0.5cm}

        \Large Questions?
    \end{center}

    \vspace{0.5cm}

    \textbf{Key Takeaways}

    \begin{itemize}
        \item Beamer is used to create presentations using LaTeX.
        \item Each slide is created using the frame environment.
        \item We can include text, equations, tables and figures.
    \end{itemize}

\end{frame}
\end{document}